\chapter{Introducción}
\label{cap:intro}

% ============================================================
\section{Motivación}
% ============================================================

En los entornos corporativos actuales, el volumen de información disponible crece de forma constante, pero no crece de forma ordenada. La documentación relevante suele estar repartida entre repositorios como Microsoft SharePoint, PDFs internos (muchos de ellos escaneados), wikis, portales de proveedores, procedimientos técnicos y fuentes públicas como el Boletín Oficial del Estado (BOE). En la práctica, el problema habitual no es ``falta de información'', sino \textbf{dispersión, formatos heterogéneos y dificultad para recuperar contenido útil} cuando se necesita.

Los métodos tradicionales de búsqueda por palabras clave tienden a fallar en situaciones comunes: cuando el usuario no conoce el término exacto (sinónimos, acrónimos internos o nomenclaturas no estandarizadas), cuando la información está contenida en documentos extensos o cuando el contenido está en PDFs escaneados y requiere OCR. Como consecuencia, tareas rutinarias ---confirmar un procedimiento, localizar un requisito normativo, recuperar una especificación o una decisión de arquitectura anterior--- se convierten en procesos lentos, repetitivos y con alto riesgo de omisión.

En paralelo, la popularización de los grandes modelos de lenguaje (LLMs, \textit{Large Language Models}) ha permitido interactuar con sistemas de información mediante lenguaje natural. Sin embargo, en un contexto corporativo, una respuesta fluida pero incorrecta tiene un coste elevado: puede inducir decisiones erróneas y generar confianza injustificada. El problema se agrava cuando el modelo no dispone de acceso al conocimiento interno, cuando la documentación cambia con frecuencia o cuando la consulta requiere precisión y respaldo documental.

Por este motivo, resulta especialmente relevante el enfoque de Generación Aumentada por Recuperación (RAG, \textit{Retrieval-Augmented Generation}), que combina la generación de lenguaje con la recuperación de información desde una base documental. En lugar de ``responder de memoria'', el sistema fundamenta la respuesta en documentos recuperados, mejorando la \textbf{utilidad}, la \textbf{trazabilidad} y reduciendo la probabilidad de \textbf{respuestas no justificadas} \cite{lewis2020rag}.

Con esta motivación, este Trabajo de Fin de Grado aborda el diseño, desarrollo e integración de un sistema RAG corporativo (\textit{Enterprise RAG System}) desplegado en infraestructura controlada (on-premise), capaz de ingerir documentos desde múltiples fuentes heterogéneas (SharePoint, BOE y web) y ofrecer a los usuarios una interfaz conversacional orientada a consultas técnicas, con recuperación semántica y soporte documental.

% ============================================================
\subsection{Problemática del Shadow AI en entornos corporativos}
\label{subsec:shadow_ai}
% ============================================================

La adopción de asistentes de IA comerciales en el ámbito laboral ha impulsado un fenómeno cada vez más habitual: el uso no autorizado de herramientas externas por parte de empleados sin supervisión del departamento de TI ni revisión legal o de cumplimiento. Este patrón se conoce como \textit{Shadow AI} y puede considerarse una extensión del \textit{Shadow IT}, con un riesgo añadido: la exposición de datos puede producirse en el mismo momento en que un usuario copia y pega información interna en un servicio externo \cite{microsoft2024worktrend}.

La magnitud del fenómeno es considerable. Según el informe \textit{Work Trend Index} de Microsoft y LinkedIn (2024), el 75\% de los trabajadores del conocimiento utiliza herramientas de IA en su trabajo, y el 78\% de ellos emplea herramientas propias no autorizadas por la organización \cite{microsoft2024worktrend}. En un análisis de Cyberhaven sobre 1,6 millones de trabajadores, se detectó que el 8,6\% de los empleados introducía datos corporativos en ChatGPT, y que el 11\% de esos datos estaba clasificado como confidencial \cite{cyberhaven2024shadow}. Además, entre marzo de 2023 y marzo de 2024, la cantidad de datos corporativos introducidos en herramientas de IA se multiplicó por 4,85, y la proporción de datos sensibles casi se triplicó, pasando del 10,7\% al 27,4\% \cite{kalisa2024shadow}.

Los riesgos principales asociados a este fenómeno son:

\begin{itemize}
    \item \textbf{Exposición de información sensible}: al introducir documentos, fragmentos de procedimientos, incidencias, datos operativos o contenido contractual en herramientas externas, la información puede salir del perímetro de seguridad corporativo. Un estudio de Komprise (2025) reveló que el 44\% de las organizaciones ha experimentado fugas de datos sensibles hacia herramientas de IA \cite{komprise2025survey}.

    \item \textbf{Riesgo regulatorio y de cumplimiento}: en sectores regulados o con datos personales/sensibles, el tratamiento de la información exige garantías sobre finalidad, acceso, ubicación y medidas de seguridad. Un estudio de PwC (2024) indicó que el 54\% de las empresas que utilizaban herramientas de IA no autorizadas recibieron escrutinio regulatorio, con multas medias de 5 millones de euros en Europa \cite{pwc2024shadowai}. El uso informal de herramientas externas introduce incertidumbre y eleva el riesgo de incumplimiento normativo, especialmente ante el RGPD y la próxima Ley de IA de la UE.

    \item \textbf{Calidad y consistencia}: los modelos generalistas pueden producir respuestas plausibles incluso cuando carecen de base documental. En ámbitos técnicos o normativos, esto se traduce en errores silenciosos con apariencia de certeza.

    \item \textbf{Ausencia de gobernanza}: sin control centralizado, la organización no sabe qué herramientas se usan, con qué datos, con qué finalidad ni con qué impacto. Según IBM, los incidentes relacionados con Shadow AI añaden un coste medio de 308.000 dólares por brecha de seguridad \cite{ibm2025databreach}.
\end{itemize}

Este proyecto aborda directamente esta problemática mediante una alternativa interna: un sistema conversacional desplegado en la infraestructura de la organización (on-premise), donde los documentos no se envían a servicios externos y donde las respuestas se apoyan en recuperación documental, reduciendo respuestas no fundamentadas y permitiendo trazabilidad hacia fuentes internas.

% ============================================================
\section{Objetivos del proyecto}
\label{sec:objetivos}
% ============================================================

El objetivo principal de este trabajo es desarrollar un asistente conversacional avanzado sustentado en una arquitectura RAG 100\% on-premise, maximizando la confidencialidad de los datos internos y reduciendo la dependencia de servicios externos.

Para alcanzar este objetivo general, se han definido los siguientes objetivos específicos:

\begin{enumerate}
    \item \textbf{Diseño de arquitectura desacoplada y modular}: desarrollar un backend en Python (FastAPI) integrado con una base de datos vectorial (Qdrant), una base de datos relacional (PostgreSQL), ejecución local de modelos (Ollama) y un componente de gestión/enrutamiento de LLMs (LiteLLM).

    \item \textbf{Ingesta multicanal y sincronización automatizada}: proporcionar capacidad de ingesta y actualización desde fuentes heterogéneas mediante Microsoft Graph para SharePoint, OCR acelerado por GPU con PaddleOCR para documentos escaneados, y extracción web semántica estática y dinámica mediante Playwright.

    \item \textbf{Acceso a normativa del BOE y preparación para MCP}: integrar la consulta del BOE dentro del sistema para permitir búsquedas y consultas normativas desde la interfaz conversacional. En la implementación actual, esta funcionalidad se realiza mediante llamadas a API. Adicionalmente, se desarrolla un servidor compatible con el enfoque de Model Context Protocol (MCP) \cite{anthropic2024mcp} como base para integraciones futuras; dicho servidor MCP funciona y puede utilizarse desde terminal con un LLM y un cliente compatible, pero no se integra todavía de forma nativa en la interfaz web utilizada.

    \item \textbf{Enrutamiento inteligente de consultas}: implementar un agente en Python (JARVIS Pipeline) capaz de detectar la intención del usuario y derivar cada consulta al modo de operación más adecuado (RAG, búsqueda web controlada, scraping, normativa, visión o conversación general), priorizando utilidad y trazabilidad.

    \item \textbf{Despliegue, operación y observabilidad}: contenerizar el sistema con Docker y habilitar métricas de uso y rendimiento mediante Prometheus y Grafana, facilitando monitorización y análisis del comportamiento del sistema.
\end{enumerate}

% ============================================================
\section{Alcance y limitaciones}
% ============================================================

\subsection{Alcance}

El proyecto cubre el ciclo completo de construcción de un sistema RAG empresarial, desde el diseño de arquitectura hasta su despliegue reproducible:

\begin{itemize}
    \item Diseño e implementación de una arquitectura basada en microservicios contenerizados.
    \item Pipeline de ingesta documental con soporte para PDFs, HTML, imágenes y documentación de SharePoint.
    \item Motor de búsqueda semántica basado en embeddings y almacenamiento vectorial.
    \item Interfaz conversacional web con autenticación corporativa (SSO).
    \item Integración con fuentes externas (SharePoint, BOE e Internet) gestionada desde el propio sistema.
    \item Stack de monitorización y observabilidad para métricas y rendimiento.
\end{itemize}

\subsection{Limitaciones}

Se identifican las siguientes limitaciones reconocidas:

\begin{itemize}
    \item \textbf{Requisitos hardware}: el sistema requiere una GPU NVIDIA con al menos 16~GB de VRAM para mantener un rendimiento aceptable en tareas de OCR e inferencia local, lo que limita su ejecución en hardware sin aceleración.

    \item \textbf{BOE: integración funcional por API, no por MCP en la interfaz web}: el sistema permite búsquedas y consultas normativas del BOE desde la interfaz conversacional mediante llamadas a API. Sin embargo, la integración MCP nativa no se incorpora en la interfaz web utilizada, debido a limitaciones actuales del cliente.

    \item \textbf{MCP validado en entorno de terminal}: el servidor MCP desarrollado se ha validado en ejecución desde terminal con un LLM y un cliente compatible, quedando su integración directa en la interfaz web como línea futura.

    \item \textbf{Sesgo en el ajuste fino}: el fine-tuning LoRA se realiza con datos sintéticos generados por el propio modelo base, lo que puede introducir sesgos y limitar la generalización.

    \item \textbf{Evaluación acotada}: la evaluación se realiza principalmente de forma cualitativa y orientada a validación funcional; benchmarks académicos formales como RAGAS \cite{es2023ragas} o ARES quedan fuera del alcance temporal del proyecto.
\end{itemize}

% ============================================================
\section{Contribuciones principales}
% ============================================================

Las contribuciones técnicas concretas de este trabajo son:

\begin{enumerate}
    \item Una arquitectura reproducible basada en microservicios, definida declarativamente en Docker Compose, con un conjunto de servicios interconectados orientados a ingesta, almacenamiento, recuperación, generación y monitorización.

    \item Un agente enrutador (JARVIS) capaz de seleccionar distintos modos de operación en función de la intención detectada, priorizando recuperación documental cuando el caso lo requiere.

    \item Un pipeline de ingesta multicanal que unifica SharePoint, scraping web, OCR y carga manual en una única base de conocimiento vectorial.

    \item Una integración funcional con el BOE dentro del sistema, permitiendo búsquedas y consultas normativas mediante API.

    \item Desarrollo y validación de un servidor MCP para consulta normativa del BOE, operativo en ejecución por terminal con un LLM y un cliente compatible, y planteado como base para integraciones MCP nativas en iteraciones futuras.

    \item Un flujo de fine-tuning con LoRA orientado a adaptar modelos base al dominio de uso definido, con las limitaciones derivadas del uso de datos sintéticos.

    \item Publicación del sistema como proyecto de código abierto en GitHub\footnote{\url{https://github.com/acidxlemons/TFG-JARVIS}}, con documentación para despliegue reproducible.
\end{enumerate}

% ============================================================
\section{Estructura de la memoria}
% ============================================================

El resto de este documento se estructura de la siguiente manera:

En el \textbf{Capítulo~\ref{cap:estado}} se expone el estado del arte de los LLMs, los sistemas RAG, los modelos de embeddings, las bases de datos vectoriales, técnicas de cuantización, LoRA y MCP.

El \textbf{Capítulo~\ref{cap:tecnologias}} describe en profundidad el marco tecnológico del proyecto, detallando cada tecnología utilizada con especial énfasis en OpenWebUI y su sistema de Pipelines, que constituye el mecanismo de extensión central del sistema.

El \textbf{Capítulo~\ref{cap:sistemas}} analiza los sistemas existentes similares al desarrollado, incluyendo asistentes comerciales (ChatGPT Enterprise, Microsoft Copilot, Google Gemini) y soluciones de código abierto (PrivateGPT, AnythingLLM, Danswer), presentando una tabla comparativa exhaustiva.

El \textbf{Capítulo~\ref{cap:funcional}} ofrece una descripción funcional completa del sistema desde la perspectiva del usuario: la interfaz conversacional, los modos de operación disponibles, escenarios de uso y el panel de administración.

El \textbf{Capítulo~\ref{cap:implementacion}} constituye el núcleo técnico del proyecto. Describe la filosofía de diseño iterativo, la arquitectura de microservicios, la distinción entre desarrollo propio y herramientas externas, el mecanismo Pipeline (hook) de OpenWebUI, el agente enrutador JARVIS, el motor RAG, las integraciones (SharePoint, BOE, web), la seguridad, el stack de observabilidad (Prometheus, Grafana), la caché Redis y el fine-tuning con LoRA.

El \textbf{Capítulo~\ref{cap:metodologia}} describe la metodología ágil seguida, las iteraciones de desarrollo, el entorno hardware, el stack tecnológico completo, la estimación de presupuesto y la estrategia de documentación.

El \textbf{Capítulo~\ref{cap:resultados}} presenta los resultados obtenidos en producción: métricas de rendimiento, evaluación funcional de cada modo de operación, análisis de alucinaciones y el impacto de la caché.

Por último, el \textbf{Capítulo~\ref{cap:conclusiones}} recoge las conclusiones, lecciones aprendidas, el grado de consecución de los objetivos y las líneas futuras de trabajo.